\section{Introdução}

O tratamento convencional de esgoto depende de comunidades microbianas para remover matéria orgânica (BOD) e converter amônia em nitrato (nitrificação). Salinidade alta -- medida por Sólidos Dissolvidos Totais (TDS) -- pode inibir esses processos \cite{schwabe2020unintended}. Em regiões como Los Angeles, medidas de conservação de água reduzem o uso interno de água nas residências, o que pode aumentar involuntariamente a salinidade do esgoto: a mesma massa de sais passa a entrar na estação de tratamento dissolvida em um volume menor de água \cite{schwabe2020unintended}.

Este trabalho analisa 15 anos (fevereiro de 2011 a março de 2026) de medições mensais de TDS, cloreto, amônia e BOD no efluente da Los Angeles--Glendale Water Reclamation Plant (LAGWRP), obtidas do portal público eSMR do California Water Boards. A pergunta inicial -- a salinidade aumentou de forma monotônica ao longo desse período? -- não se sustentou diante dos próprios dados (Seção de Resultados), e foi substituída pela pergunta que efetivamente orienta este trabalho: \textbf{o padrão de salinidade observado é explicado por ciclos de seca da Califórnia via diluição (menos vazão diluindo a mesma carga de sais), e essa tendência está associada a sinais mensuráveis de perda de desempenho biológico do tratamento?}

Este trabalho contribui com: (1) uma bateria de 35 métodos complementares -- estatística de tendência, séries temporais, aprendizado de máquina, e nove métodos desenhados especificamente para a pergunta causal (WRTDS, balanço de massa, cenários climáticos, entre outros) -- em vez de um único modelo; (2) confirmação, com uma variável externa independente (PDSI do NOAA), de que a série segue um padrão \textbf{cíclico por regime} ligado a secas, não uma tendência monotônica, e confirmação do mecanismo causal (diluição) via balanço de massa não-circular; (3) previsão de TDS para +10/+15/+20 anos como faixa condicionada a cenários climáticos, em vez de valor único; (4) correlação entre TDS e dois indicadores de desempenho biológico (amônia e BOD), com e sem remoção de tendência comum; (5) interpretação à luz da literatura sobre conservação de água \cite{schwabe2020unintended} e sobre o papel relativo da água de origem \cite{scsc2018tds}.
